\documentclass[conference]{IEEEtran}
\IEEEoverridecommandlockouts
% The preceding line is only needed to identify funding in the first footnote. If that is unneeded, please comment it out.
\usepackage{amsmath,amssymb,amsfonts}
\usepackage{algorithmic}
\usepackage{graphicx}
\usepackage{textcomp}
% \usepackage{cite}
\usepackage{biblatex}
\usepackage{subcaption}
\usepackage{hyperref}
\usepackage{xcolor}
\def\BibTeX{{\rm B\kern-.05em{\sc i\kern-.025em b}\kern-.08em
    T\kern-.1667em\lower.7ex\hbox{E}\kern-.125emX}}

    
\begin{document}

\title{Title}

\author{
    \IEEEauthorblockN{Luka Antonyshyn, Anthony Beca, Aoran Jiao, and Elliot Preston-Krebs}
    \IEEEauthorblockA{\textit{Institute for Aerospace Studies, University of Toronto}, Toronto, Canada \\
    \{luka.antonyshyn, anthony.beca, aoran.jiao, elliot.prestonkrebs\}@robotics.utias.utoronto.ca}
}

\maketitle

\begin{abstract}
\end{abstract}

\begin{IEEEkeywords}
\end{IEEEkeywords}

\section{Introduction}
% give a high level overview of your project - is it a technique, or the application of a specific algorithm, etc.? Who (worldwide) the major players in the area? Why is the topic interesting?
Adversarial games in continuous spaces have long been studied under the field of differential game theory, developed by Rufus Isaac's in the 1950s~\cite{}.
These games, contrary to their discrete classical game theoretical counterpoints, have the evolution of the game evolve according to differential equations.
Differential games have been used in several application spaces, such as perimeter defense, risk-averse landing path plannign for airplanes, and ... \cite{}.
In robotics, differential games find many applications in military ...

Recently, more interest has been taken in multi-agent differential games~\cite{}. 
These, of course, model more realistic scenarios and engagements. 
However, they must often settle for sub-optimal solutions due to the curse of dimensionality and the difficulty of somputing Hamilton-Jacobi-Isaacs solutions.~\cite{}

One such multi-agent game is the multi-pursuer vision-based pursuit evasion game.
 In this formulation, the evader, or evaders, are often treated as abstract entities that can move at arbitrary speeds. 
 This forces the pursuing team to plan for the worst-case scenario.
 As such, the problem devolves into a type of constrained frontier expansion problem on a known path.
 This makes the formulation of interest to motion planning researchers.
 
Many solutions have been proposed for this problem, from learning-based methods, sampling-based approaches, and graph-search algorithms~\cite{}.
However, in the literature, practitioners tend to compare only against solutions they have readily available - ignoring suitable baselines from other problem formulations and only evaluating on limited subsets of the total problem space.
As such, the state of the art is lacking in honest analysis and comparison of different approaches to multi-agent differential games.

Following the growth of popularity of game theory, an effort to evaluate the performance of different strategies was undertaken on Axelrod's tournament - a contest in which different policies were pitted against each other in varying conditions to compare outcomes.
Here, we pursue a simplified version of Axelrod's tournament on one differential game.
We implement four approaches to the multi-pursuer observability-based pursuit evasion game, and compare them on numerous scenarios using common metrics.
We identify and analyse the behaviours of each method under various conditions that challenge their formulations, and use the information to recommend steps for future works.

The remainder of this work is organized as follows; in Section \ref{}

\section{Background}
% Describe the necessary background of your topic. How did researchers in the field derive the current approach(es)? Explain the relevance of the topic to motion planning

\subsection{Pursuit-Evasion Games}



\section{Literature Review}
% Provide a structured review of seminal and recent work in your topic area. You should attemp to provide a (brief) taxonomy of research results, organized in such a way that someone reading your report will be able ot understand some common research threads
\subsection{Sampling-based Approaches}

\subsection{Graph-based Approaches}

\subsection{Learning-based Approaches}

\section{Implementation}
% What did you implement, and what were the results you obtained? What information is presented in the tables, figures, and diagrams you have included? This is the 'meat' of the report - when writing this section, step back and consider how best to convey your results to someone who is new to the topic.

\section{Open Problems}
% Based upon your work, what are some (specific) open problems in the topic area? Are there any major obstacles to progress (e.g. insufficient compute, unsolved theoretical issues)? Can you suggest possible ways to move forward?

\printbibliography


\end{document}

